\documentclass[11pt,a4paper]{article}
\usepackage[utf8]{inputenc}
\usepackage[T1]{fontenc}
\usepackage[spanish,es-minimal,es-noshorthands]{babel}
\AtBeginDocument{\renewcommand{\sin}{\operatorname{sen}}\renewcommand{\arcsin}{\operatorname{arcsen}}}
\usepackage{lmodern}
\usepackage[sfdefault,lining]{FiraSans}
\usepackage[a4paper,top=2.3cm,bottom=2.2cm,left=1.8cm,right=1.8cm,headheight=1.3cm,headsep=0.4cm,footskip=1cm]{geometry}
\usepackage{amsmath,amssymb}
\usepackage{icomma}
\usepackage{graphicx}
\usepackage[table]{xcolor}
\usepackage{tikz}
\usepackage{fontawesome5}
\usepackage{ccicons}
\usepackage{fancyhdr}
\usepackage{lastpage}
\usepackage{tasks}
\usepackage{needspace}
\usepackage[most]{tcolorbox}
\usepackage{array}
\usepackage[hidelinks]{hyperref}
\AtBeginDocument{\hypersetup{pdftitle={Ficha de ejercicios de la unidad 11: Probabilidad. 4º ESO},pdfsubject={Matemáticas A 4º ESO},pdfkeywords={probabilidad, experimento aleatorio, espacio muestral, sucesos, unión, intersección, suceso contrario, frecuencia relativa, ley de Laplace, experiencias compuestas, experiencias independientes, experiencias dependientes, probabilidad condicionada, diagrama de árbol, tabla de contingencia}}}

\definecolor{navy}{HTML}{1B2748}
\definecolor{accent}{HTML}{E8603A}
\definecolor{teal}{HTML}{1F9E90}
\definecolor{soft}{HTML}{EEF1F8}
\definecolor{muted}{HTML}{5D6685}

\setlength{\parindent}{0pt}
\setlength{\parskip}{3pt}
\settasks{label=\alph*), label-format=\bfseries\color{navy}, label-width=1.3em, item-indent=1.8em, column-sep=0.6em, after-item-skip=5pt, before-skip=3pt, after-skip=4pt}

% ---------- capçalera i peu ----------
\newcommand{\llicencia}{\href{https://creativecommons.org/licenses/by-nc-sa/4.0/deed.es}{\ccbyncsa}\ \ Licencia Creative Commons Atribución-NoComercial-CompartirIgual 4.0 Internacional (CC BY-NC-SA 4.0)}
\pagestyle{fancy}
\fancyhf{}
\fancyhead[L]{\raisebox{-0.35\height}{\includegraphics[height=0.95cm]{logo-gva.png}}}
\fancyhead[C]{\small\color{muted}Ficha de ejercicios · Unidad 11: Probabilidad · 4\textsuperscript{o} ESO}
\fancyhead[R]{\raisebox{-0.35\height}{\includegraphics[height=0.95cm]{logo-vermellar.png}}}
\fancyfoot[L]{\scriptsize\color{muted}\llicencia}
\fancyfoot[R]{\small\color{muted}\thepage}
\renewcommand{\headrulewidth}{0.4pt}
\renewcommand{\headrule}{\color{muted!40}\hrule height \headrulewidth}
\fancypagestyle{primera}{\fancyhf{}\renewcommand{\headrulewidth}{0pt}
  \fancyfoot[L]{\scriptsize\color{muted}\llicencia}
  \fancyfoot[R]{\small\color{muted}\thepage}}

% ---------- nivells i exercicis ----------
\newcommand{\estrelles}[1]{{\color{accent}\ifcase#1\or\faStar\faStar[regular]\faStar[regular]\or\faStar\faStar\faStar[regular]\or\faStar\faStar\faStar\fi}}
\newcounter{exer}
\newcommand{\ex}[1]{\par\needspace{4\baselineskip}\medskip\stepcounter{exer}\noindent{\color{navy}\bfseries\theexer.}\ \ {\small\estrelles{#1}}\ \ }

\newcommand{\seccio}[2]{\par\needspace{10\baselineskip}\bigskip
  \begin{tcolorbox}[colback=navy,colframe=navy,coltext=white,arc=2mm,boxrule=0pt,left=3mm,right=3mm,top=1.2mm,bottom=1.2mm]
  {\large\bfseries #1\quad #2}
  \end{tcolorbox}\vspace{-2pt}}

\newcolumntype{C}[1]{>{\centering\arraybackslash}p{#1}}
\usepackage{pifont}
\usetikzlibrary{trees}
% quadre de recordatori (només a la fitxa d'exercicis)
\newcommand{\recorda}[1]{\par\smallskip\begin{tcolorbox}[colback=accent!7,colframe=accent!55,arc=2mm,boxrule=0.5pt,left=2.5mm,right=2.5mm,top=1mm,bottom=1mm]\small{\color{accent}\faLightbulb[regular]}\ \textbf{Recuerda:}\ #1\end{tcolorbox}}
\newcommand{\res}[1]{#1}\newcommand{\rest}[1]{#1}

\def\fmneed{14}
% =====================================================================
% Capçalera FentMates de les fitxes: un color per a cada curs.
% Els colors són els de la pàgina inicial (LEVELS a assets/js/materials.js);
% el logo ve de _eines/logo/pdf/ (python3 _eines/logo/genera.py).
%
% Ús, al final del preàmbul de la unitat:
%   % =====================================================================
% Capçalera FentMates de les fitxes: un color per a cada curs.
% Els colors són els de la pàgina inicial (LEVELS a assets/js/materials.js);
% el logo ve de _eines/logo/pdf/ (python3 _eines/logo/genera.py).
%
% Ús, al final del preàmbul de la unitat:
%   % =====================================================================
% Capçalera FentMates de les fitxes: un color per a cada curs.
% Els colors són els de la pàgina inicial (LEVELS a assets/js/materials.js);
% el logo ve de _eines/logo/pdf/ (python3 _eines/logo/genera.py).
%
% Ús, al final del preàmbul de la unitat:
%   \input{../comu/fentmates.tex}
%   \fmcurs{1eso}
% i, al començament del document:
%   \capcalera{Matemàtiques}{1r ESO}{Fitxa d'exercicis · Unitat 2}{Divisibilitat}
% Opcional abans de \input: \def\fmneed{14} (espai mínim abans de cada secció).
% =====================================================================
\makeatletter
\def\fm@logodir{../../logo/pdf}
\providecommand{\fmneed}{10}
\def\fm@cursnom{}

% clau del curs, color i color del text damunt del color
\newcommand{\fm@curs}[3]{\@namedef{fm@color@#1}{#2}\@namedef{fm@text@#1}{#3}}
\fm@curs{1eso}{E6B84A}{navy}
\fm@curs{2eso}{3B82F6}{white}
\fm@curs{3eso}{FF7A50}{white}
\fm@curs{4eso-a}{A78BFA}{white}
\fm@curs{4eso-b}{6D28D9}{white}
\fm@curs{1batx-mi}{9B1C1C}{white}
\fm@curs{1batx-ccss}{EF6B6B}{white}
\fm@curs{2batx-mii}{1E6B47}{white}
\fm@curs{2batx-ccss}{2BB3A3}{white}

\newcommand{\fmcurs}[1]{%
  \@ifundefined{fm@color@#1}{\PackageError{fentmates}{Curs desconegut: #1}{}}{}%
  \edef\fm@tmp{\noexpand\definecolor{curs}{HTML}{\@nameuse{fm@color@#1}}%
               \noexpand\colorlet{curstext}{\@nameuse{fm@text@#1}}}\fm@tmp
  \colorlet{cursfosc}{curs!72!black}%
  \def\fm@logo{\fm@logodir/vertical-#1.pdf}%
  \def\fm@marca{\fm@logodir/marca-#1.pdf}}

% ---------- capçalera de la primera pàgina ----------
% #1 matèria (petit), #2 curs, #3 tipus de document i unitat, #4 títol
\newcommand{\capcalera}[4]{%
  \gdef\fm@cursnom{#2}%
  \par\noindent
  \begin{tcolorbox}[colback=curs,colframe=curs,boxrule=0pt,arc=2.2mm,width=3.5cm,height=2.15cm,
      valign=center,left=3mm,right=2mm,top=1mm,bottom=1mm,nobeforeafter,box align=center]
    \color{curstext}{\footnotesize #1}\par\vspace{1pt}{\fontsize{19}{21}\selectfont\bfseries #2\par}%
  \end{tcolorbox}\hspace{5mm}%
  \begin{minipage}[c]{\dimexpr\linewidth-3.5cm-5mm-2.9cm\relax}
    {\color{cursfosc}\small\bfseries\MakeUppercase{#3}}\par\vspace{2pt}
    {\color{navy}\fontsize{20}{23}\selectfont\bfseries #4\par}
  \end{minipage}\hfill
  \begin{minipage}[c]{2.6cm}\raggedleft\includegraphics[width=2.6cm]{\fm@logo}\end{minipage}\par
  \vspace{2pt}}

% ---------- barres de secció del color del curs ----------
\renewcommand{\seccio}[2]{\par\needspace{\fmneed\baselineskip}\bigskip
  \begin{tcolorbox}[colback=curs,colframe=curs,coltext=curstext,arc=2mm,boxrule=0pt,left=3mm,right=3mm,top=1.2mm,bottom=1.2mm]
  {\large\bfseries #1\quad #2}
  \end{tcolorbox}\vspace{-2pt}}

% ---------- estrelles de dificultat del color del curs (a Batxillerat no se'n posen) ----------
\@ifundefined{estrelles}{}{%
  \renewcommand{\estrelles}[1]{{\color{curs}\ifcase#1\or\faStar\faStar[regular]\faStar[regular]\or\faStar\faStar\faStar[regular]\or\faStar\faStar\faStar\fi}}}

% ---------- capçalera i peu de les pàgines ----------
% primera pàgina: logos institucionals al peu
\fancypagestyle{primera}{\fancyhf{}\renewcommand{\headrulewidth}{0pt}
  \fancyfoot[L]{\raisebox{-0.32\height}{\includegraphics[height=0.78cm]{logo-gva.png}}\hspace{5mm}%
    \raisebox{-0.32\height}{\includegraphics[height=0.78cm]{logo-vermellar.png}}}
  \fancyfoot[R]{\scriptsize\color{muted}\href{https://creativecommons.org/licenses/by-nc-sa/4.0/deed.ca}{\ccbyncsa}\ \ CC BY-NC-SA 4.0 · fentmates.org\qquad\small\thepage}}
% la resta: el símbol a l'esquerra i el curs a la dreta, amb el títol al centre com abans
\fancyhead[L]{\raisebox{-0.3\height}{\includegraphics[height=0.62cm]{\fm@marca}}}
\fancyhead[R]{\small\bfseries\color{cursfosc}\fm@cursnom}
\renewcommand{\headrule}{\color{curs!45}\hrule height \headrulewidth}
\makeatother

%   \fmcurs{1eso}
% i, al començament del document:
%   \capcalera{Matemàtiques}{1r ESO}{Fitxa d'exercicis · Unitat 2}{Divisibilitat}
% Opcional abans de \input: \def\fmneed{14} (espai mínim abans de cada secció).
% =====================================================================
\makeatletter
\def\fm@logodir{../../logo/pdf}
\providecommand{\fmneed}{10}
\def\fm@cursnom{}

% clau del curs, color i color del text damunt del color
\newcommand{\fm@curs}[3]{\@namedef{fm@color@#1}{#2}\@namedef{fm@text@#1}{#3}}
\fm@curs{1eso}{E6B84A}{navy}
\fm@curs{2eso}{3B82F6}{white}
\fm@curs{3eso}{FF7A50}{white}
\fm@curs{4eso-a}{A78BFA}{white}
\fm@curs{4eso-b}{6D28D9}{white}
\fm@curs{1batx-mi}{9B1C1C}{white}
\fm@curs{1batx-ccss}{EF6B6B}{white}
\fm@curs{2batx-mii}{1E6B47}{white}
\fm@curs{2batx-ccss}{2BB3A3}{white}

\newcommand{\fmcurs}[1]{%
  \@ifundefined{fm@color@#1}{\PackageError{fentmates}{Curs desconegut: #1}{}}{}%
  \edef\fm@tmp{\noexpand\definecolor{curs}{HTML}{\@nameuse{fm@color@#1}}%
               \noexpand\colorlet{curstext}{\@nameuse{fm@text@#1}}}\fm@tmp
  \colorlet{cursfosc}{curs!72!black}%
  \def\fm@logo{\fm@logodir/vertical-#1.pdf}%
  \def\fm@marca{\fm@logodir/marca-#1.pdf}}

% ---------- capçalera de la primera pàgina ----------
% #1 matèria (petit), #2 curs, #3 tipus de document i unitat, #4 títol
\newcommand{\capcalera}[4]{%
  \gdef\fm@cursnom{#2}%
  \par\noindent
  \begin{tcolorbox}[colback=curs,colframe=curs,boxrule=0pt,arc=2.2mm,width=3.5cm,height=2.15cm,
      valign=center,left=3mm,right=2mm,top=1mm,bottom=1mm,nobeforeafter,box align=center]
    \color{curstext}{\footnotesize #1}\par\vspace{1pt}{\fontsize{19}{21}\selectfont\bfseries #2\par}%
  \end{tcolorbox}\hspace{5mm}%
  \begin{minipage}[c]{\dimexpr\linewidth-3.5cm-5mm-2.9cm\relax}
    {\color{cursfosc}\small\bfseries\MakeUppercase{#3}}\par\vspace{2pt}
    {\color{navy}\fontsize{20}{23}\selectfont\bfseries #4\par}
  \end{minipage}\hfill
  \begin{minipage}[c]{2.6cm}\raggedleft\includegraphics[width=2.6cm]{\fm@logo}\end{minipage}\par
  \vspace{2pt}}

% ---------- barres de secció del color del curs ----------
\renewcommand{\seccio}[2]{\par\needspace{\fmneed\baselineskip}\bigskip
  \begin{tcolorbox}[colback=curs,colframe=curs,coltext=curstext,arc=2mm,boxrule=0pt,left=3mm,right=3mm,top=1.2mm,bottom=1.2mm]
  {\large\bfseries #1\quad #2}
  \end{tcolorbox}\vspace{-2pt}}

% ---------- estrelles de dificultat del color del curs (a Batxillerat no se'n posen) ----------
\@ifundefined{estrelles}{}{%
  \renewcommand{\estrelles}[1]{{\color{curs}\ifcase#1\or\faStar\faStar[regular]\faStar[regular]\or\faStar\faStar\faStar[regular]\or\faStar\faStar\faStar\fi}}}

% ---------- capçalera i peu de les pàgines ----------
% primera pàgina: logos institucionals al peu
\fancypagestyle{primera}{\fancyhf{}\renewcommand{\headrulewidth}{0pt}
  \fancyfoot[L]{\raisebox{-0.32\height}{\includegraphics[height=0.78cm]{logo-gva.png}}\hspace{5mm}%
    \raisebox{-0.32\height}{\includegraphics[height=0.78cm]{logo-vermellar.png}}}
  \fancyfoot[R]{\scriptsize\color{muted}\href{https://creativecommons.org/licenses/by-nc-sa/4.0/deed.ca}{\ccbyncsa}\ \ CC BY-NC-SA 4.0 · fentmates.org\qquad\small\thepage}}
% la resta: el símbol a l'esquerra i el curs a la dreta, amb el títol al centre com abans
\fancyhead[L]{\raisebox{-0.3\height}{\includegraphics[height=0.62cm]{\fm@marca}}}
\fancyhead[R]{\small\bfseries\color{cursfosc}\fm@cursnom}
\renewcommand{\headrule}{\color{curs!45}\hrule height \headrulewidth}
\makeatother

%   \fmcurs{1eso}
% i, al començament del document:
%   \capcalera{Matemàtiques}{1r ESO}{Fitxa d'exercicis · Unitat 2}{Divisibilitat}
% Opcional abans de \input: \def\fmneed{14} (espai mínim abans de cada secció).
% =====================================================================
\makeatletter
\def\fm@logodir{../../logo/pdf}
\providecommand{\fmneed}{10}
\def\fm@cursnom{}

% clau del curs, color i color del text damunt del color
\newcommand{\fm@curs}[3]{\@namedef{fm@color@#1}{#2}\@namedef{fm@text@#1}{#3}}
\fm@curs{1eso}{E6B84A}{navy}
\fm@curs{2eso}{3B82F6}{white}
\fm@curs{3eso}{FF7A50}{white}
\fm@curs{4eso-a}{A78BFA}{white}
\fm@curs{4eso-b}{6D28D9}{white}
\fm@curs{1batx-mi}{9B1C1C}{white}
\fm@curs{1batx-ccss}{EF6B6B}{white}
\fm@curs{2batx-mii}{1E6B47}{white}
\fm@curs{2batx-ccss}{2BB3A3}{white}

\newcommand{\fmcurs}[1]{%
  \@ifundefined{fm@color@#1}{\PackageError{fentmates}{Curs desconegut: #1}{}}{}%
  \edef\fm@tmp{\noexpand\definecolor{curs}{HTML}{\@nameuse{fm@color@#1}}%
               \noexpand\colorlet{curstext}{\@nameuse{fm@text@#1}}}\fm@tmp
  \colorlet{cursfosc}{curs!72!black}%
  \def\fm@logo{\fm@logodir/vertical-#1.pdf}%
  \def\fm@marca{\fm@logodir/marca-#1.pdf}}

% ---------- capçalera de la primera pàgina ----------
% #1 matèria (petit), #2 curs, #3 tipus de document i unitat, #4 títol
\newcommand{\capcalera}[4]{%
  \gdef\fm@cursnom{#2}%
  \par\noindent
  \begin{tcolorbox}[colback=curs,colframe=curs,boxrule=0pt,arc=2.2mm,width=3.5cm,height=2.15cm,
      valign=center,left=3mm,right=2mm,top=1mm,bottom=1mm,nobeforeafter,box align=center]
    \color{curstext}{\footnotesize #1}\par\vspace{1pt}{\fontsize{19}{21}\selectfont\bfseries #2\par}%
  \end{tcolorbox}\hspace{5mm}%
  \begin{minipage}[c]{\dimexpr\linewidth-3.5cm-5mm-2.9cm\relax}
    {\color{cursfosc}\small\bfseries\MakeUppercase{#3}}\par\vspace{2pt}
    {\color{navy}\fontsize{20}{23}\selectfont\bfseries #4\par}
  \end{minipage}\hfill
  \begin{minipage}[c]{2.6cm}\raggedleft\includegraphics[width=2.6cm]{\fm@logo}\end{minipage}\par
  \vspace{2pt}}

% ---------- barres de secció del color del curs ----------
\renewcommand{\seccio}[2]{\par\needspace{\fmneed\baselineskip}\bigskip
  \begin{tcolorbox}[colback=curs,colframe=curs,coltext=curstext,arc=2mm,boxrule=0pt,left=3mm,right=3mm,top=1.2mm,bottom=1.2mm]
  {\large\bfseries #1\quad #2}
  \end{tcolorbox}\vspace{-2pt}}

% ---------- estrelles de dificultat del color del curs (a Batxillerat no se'n posen) ----------
\@ifundefined{estrelles}{}{%
  \renewcommand{\estrelles}[1]{{\color{curs}\ifcase#1\or\faStar\faStar[regular]\faStar[regular]\or\faStar\faStar\faStar[regular]\or\faStar\faStar\faStar\fi}}}

% ---------- capçalera i peu de les pàgines ----------
% primera pàgina: logos institucionals al peu
\fancypagestyle{primera}{\fancyhf{}\renewcommand{\headrulewidth}{0pt}
  \fancyfoot[L]{\raisebox{-0.32\height}{\includegraphics[height=0.78cm]{logo-gva.png}}\hspace{5mm}%
    \raisebox{-0.32\height}{\includegraphics[height=0.78cm]{logo-vermellar.png}}}
  \fancyfoot[R]{\scriptsize\color{muted}\href{https://creativecommons.org/licenses/by-nc-sa/4.0/deed.ca}{\ccbyncsa}\ \ CC BY-NC-SA 4.0 · fentmates.org\qquad\small\thepage}}
% la resta: el símbol a l'esquerra i el curs a la dreta, amb el títol al centre com abans
\fancyhead[L]{\raisebox{-0.3\height}{\includegraphics[height=0.62cm]{\fm@marca}}}
\fancyhead[R]{\small\bfseries\color{cursfosc}\fm@cursnom}
\renewcommand{\headrule}{\color{curs!45}\hrule height \headrulewidth}
\makeatother

\fmcurs{4eso-a}

\begin{document}
\thispagestyle{primera}

\capcalera{Matemáticas A}{4\textsuperscript{o} ESO}{Ficha de ejercicios · Unidad 11}{Probabilidad}
\vspace{10pt}
Nombre y apellidos: \rule{6.6cm}{0.4pt}\hfill Grupo: \rule{1.6cm}{0.4pt}\hfill Fecha: \rule{2.2cm}{0.4pt}
\vspace{8pt}

\begin{tcolorbox}[colback=soft,colframe=soft,arc=2mm,left=3mm,right=3mm,top=1.5mm,bottom=1.5mm]
\small Ejercicios graduados por niveles de dificultad:\quad \estrelles{1}\ básico \qquad \estrelles{2}\ medio \qquad \estrelles{3}\ avanzado.\\[2pt] Resuélvelos en el cuaderno mostrando todos los pasos: escribe el espacio muestral, haz un diagrama de árbol o una tabla cuando haga falta y da las probabilidades en forma de fracción irreducible (o en decimal, si los datos lo son). Las secciones siguen el orden de los apuntes de la unidad (11.1--11.7).
\end{tcolorbox}

% =====================================================
\seccio{1}{Experimentos}

\ex{1} Indica si son experimentos aleatorios o deterministas:\begin{tasks}(2)
  \task Lanzar un dado y anotar el resultado.
  \task Medir la temperatura a la que hierve el agua al nivel del mar.
  \task Sacar una carta de una baraja sin mirar.
  \task Calcular cuánto pagaremos por 3 kg de manzanas a 2 €/kg.
  \task Adivinar el número ganador de la lotería.
  \task Dejar caer una pelota y mirar si cae al suelo.
\end{tasks}\recorda{Un experimento es \textbf{aleatorio} si no podemos predecir su resultado, aunque lo repitamos en las mismas condiciones. Es \textbf{determinista} si el resultado es siempre el mismo y lo podemos saber antes de hacerlo.}

\ex{1} En una bolsa hay 9 bolas rojas y 1 azul. Sacamos una bola sin mirar. Describe cada suceso con una de estas palabras: \emph{seguro}, \emph{muy probable}, \emph{poco probable}, \emph{imposible}.\par\smallskip\begin{center}\begin{tikzpicture}[x=0.42cm,y=0.42cm,baseline=(current bounding box.center)]
\draw[muted,thick,fill=soft] (-0.2,0.1) -- (-0.2,-1.70) to[out=-90,in=180] (2.50,-2.30) to[out=0,in=-90] (5.20,-1.70) -- (5.20,0.1);
\draw[muted,line width=1.6pt,line cap=round] (-0.4,0.1) -- (5.40,0.1);
\filldraw[draw=black!50,fill=red!75!black] (0.50,-0.45) circle (0.38);
\filldraw[draw=black!50,fill=red!75!black] (1.50,-0.45) circle (0.38);
\filldraw[draw=black!50,fill=red!75!black] (2.50,-0.45) circle (0.38);
\filldraw[draw=black!50,fill=red!75!black] (3.50,-0.45) circle (0.38);
\filldraw[draw=black!50,fill=red!75!black] (4.50,-0.45) circle (0.38);
\filldraw[draw=black!50,fill=red!75!black] (0.50,-1.45) circle (0.38);
\filldraw[draw=black!50,fill=red!75!black] (1.50,-1.45) circle (0.38);
\filldraw[draw=black!50,fill=red!75!black] (2.50,-1.45) circle (0.38);
\filldraw[draw=black!50,fill=red!75!black] (3.50,-1.45) circle (0.38);
\filldraw[draw=black!50,fill=blue!70!black] (4.50,-1.45) circle (0.38);
\end{tikzpicture}\end{center}\begin{tasks}(4)
  \task Sacar roja.
  \task Sacar azul.
  \task Sacar verde.
  \task Sacar roja o azul.
\end{tasks}

\ex{1} Escribe el espacio muestral de estos experimentos:\begin{tasks}(2)
  \task Lanzar una moneda.
  \task Lanzar un dado de cuatro caras numeradas del 1 al 4.
  \task Girar una ruleta con 5 sectores iguales: rojo, azul, verde, amarillo y negro.
  \task Sacar una bola de una bolsa con bolas rojas y azules.
  \task Elegir al azar un día de la semana.
  \task Preguntar a alguien cuántas hermanas tiene (como máximo, 3).
\end{tasks}\recorda{El \textbf{espacio muestral} $E$ es el conjunto de todos los resultados posibles, y un \textbf{suceso} es un subconjunto de $E$ (un grupo de resultados).}

\ex{2} Marc dice que lanzar un dado es un experimento determinista, «porque ya sabemos que saldrá un número del 1 al 6». ¿Tiene razón? Explícalo.

\ex{2} Lanzamos dos monedas. Haz un diagrama de árbol y escribe el espacio muestral y los sucesos $A=$ «al menos una cara», $B=$ «ninguna cara» y $C=$ «las dos iguales».

\ex{2} Una familia tiene tres hijos. Anotamos si cada uno es niña (A) o niño (O), por orden de nacimiento. Haz un diagrama de árbol, escribe el espacio muestral y el suceso «exactamente dos niñas».

\ex{2} Un menú tiene dos primeros platos, sopa (S) o ensalada (E), y tres segundos, carne (C), pescado (P) o tortilla (T). Elegimos el menú al azar. Haz un diagrama de árbol, escribe el espacio muestral y el suceso «el menú lleva pescado».

\ex{3} Formamos números de tres cifras distintas con las cifras 1, 2 y 3. Escribe el espacio muestral y los sucesos «el número es par» y «el número es mayor que 200».


% =====================================================
\seccio{2}{Sucesos}

\ex{1} Sacamos una bola de una urna con 12 bolas numeradas del 1 al 12. Escribe estos sucesos y di de qué tipo es cada uno (elemental, compuesto, seguro o imposible):\begin{tasks}(3)
  \task $A=$ «par»
  \task $B=$ «múltiplo de 5»
  \task $C=$ «mayor que 12»
  \task $D=$ «divisor de 12»
  \task $F=$ «menor que 13»
  \task $G=$ «múltiplo de 11»
\end{tasks}

\ex{1} Lanzamos un dado de 8 caras numeradas del 1 al 8. Escribe el espacio muestral y los sucesos:\begin{tasks}(4)
  \task «par»
  \task «primo»
  \task «mayor que 5»
  \task «divisor de 8»
\end{tasks}

\ex{1} Lanzamos un dado. Sean $A=$ «múltiplo de 3» y $B=$ «menor que 3». Calcula $A\cup B$, $A\cap B$, $\overline{A}$ y $\overline{B}$. ¿Son incompatibles? ¿Son contrarios?\recorda{$A\cup B$ («$A$ o $B$»): elementos de $A$ o de $B$. \ $A\cap B$ («$A$ y $B$»): elementos comunes. \ $\overline{A}$ («no $A$»): elementos de $E$ que no están en $A$. Si $A\cap B=\varnothing$, son \textbf{incompatibles}.}

\ex{1} En una urna hay diez bolas numeradas del 1 al 10. Sean $A=$ «par» y $B=$ «primo». Calcula $A\cup B$, $A\cap B$, $\overline{A}$ y $\overline{B}$.

\ex{2} Sacamos una carta de una baraja española de 40 cartas. Sean $O=$ «oros» y $A=$ «as». ¿Cuántas cartas tiene cada uno de estos sucesos: $O\cap A$, $O\cup A$ y $\overline{O}$?

\ex{2} Lanzamos tres monedas. Sean $A=$ «al menos dos caras» y $B=$ «la primera es cruz». Escribe $A$, $B$, $A\cap B$, $A\cup B$ y $\overline{A}$.

\ex{3} Sacamos una bola de una urna con bolas del 1 al 12. Sean $A=$ «múltiplo de 2» y $B=$ «múltiplo de 3». Escribe con operaciones y calcula: a) «es múltiplo de 2 pero no de 3»; b) «no es múltiplo ni de 2 ni de 3». Comprueba que b) es el contrario de $A\cup B$.


% =====================================================
\seccio{3}{Probabilidad de un suceso}

\ex{1} Hemos lanzado un dado 50 veces y hemos obtenido estos resultados. Completa la tabla con las frecuencias relativas (en decimal) y calcula la frecuencia relativa de «sacar par».\par\smallskip\begin{center}{\small\renewcommand{\arraystretch}{1.2}\begin{tabular}{l|cccccc}
\rowcolor{soft} Resultado & 1 & 2 & 3 & 4 & 5 & 6\\\hline
Frec. absoluta & 7 & 9 & 8 & 10 & 6 & 10\\
Frec. relativa &  &  &  &  &  & \\
\end{tabular}}\end{center}\recorda{Si un suceso ocurre $f$ veces en $n$ repeticiones, su \textbf{frecuencia relativa} es $h=\frac{f}{n}$.}

\ex{1} ¿Cuáles de estos números pueden ser la probabilidad de un suceso?\begin{tasks}(4)
  \task $0{,}7$
  \task $1{,}2$
  \task $-0{,}1$
  \task $\frac32$
  \task $0$
  \task $100\,\%$
  \task $\frac54$
  \task $\frac{99}{100}$
\end{tasks}

\ex{1} Calcula la probabilidad del suceso contrario:\begin{tasks}(3)
  \task $P(A)=0{,}62$
  \task $P(B)=\frac37$
  \task $P(C)=15\,\%$
\end{tasks}\recorda{$P(\overline{A})=1-P(A)$ \qquad $P(A\cup B)=P(A)+P(B)-P(A\cap B)$ \qquad Si son incompatibles, $P(A\cup B)=P(A)+P(B)$.}

\ex{2} Hemos lanzado una chincheta varias veces y hemos anotado cuántas ha caído con la punta hacia arriba. Completa la tabla y estima la probabilidad de que caiga con la punta hacia arriba.\par\smallskip\begin{center}{\small\renewcommand{\arraystretch}{1.2}\begin{tabular}{l|cccc}
\rowcolor{soft} Lanzamientos & 20 & 100 & 500 & 1000\\\hline
Punta arriba & 13 & 58 & 297 & 601\\
Frec. relativa &  &  &  & \\
\end{tabular}}\end{center}\recorda{\textbf{Ley de los grandes números}: cuando repetimos un experimento muchas veces, la frecuencia relativa de un suceso se estabiliza alrededor de su probabilidad.}

\ex{2} Un jugador de balonmano ha marcado 45 de los 60 penaltis que ha lanzado. Estima la probabilidad de que marque un penalti. ¿Cuántos penaltis esperas que marque de los 20 próximos?

\ex{2} Un dado está trucado. Las probabilidades del 1 al 5 son 0,15, 0,1, 0,2, 0,15 y 0,1. Calcula $P(6)$ y la probabilidad de sacar par.

\ex{2} Sean $A$ y $B$ dos sucesos con $P(A)=0{,}4$, $P(B)=0{,}5$ y $P(A\cap B)=0{,}15$. Calcula $P(A\cup B)$, la probabilidad de que no ocurra ninguno de los dos y la probabilidad de que ocurra $A$ pero no $B$.

\ex{2} a) $A$ y $B$ son incompatibles, con $P(A)=0{,}3$ y $P(B)=0{,}45$. Calcula $P(A\cup B)$. \ b) ¿Pueden ser incompatibles dos sucesos con $P(A)=0{,}6$ y $P(B)=0{,}5$? ¿Por qué?

\ex{3} En una bolsa hay 30 bolas, rojas y blancas. Sacamos una bola, anotamos su color y la devolvemos a la bolsa. Después de 300 extracciones han salido 120 rojas. ¿Cuántas bolas rojas debe de haber?

\ex{3} Sabemos que $P(A\cup B)=0{,}9$, $P(\overline{A})=0{,}3$ y $P(\overline{B})=0{,}4$. Calcula $P(A)$, $P(B)$ y $P(A\cap B)$. ¿Son incompatibles?


% =====================================================
\seccio{4}{Cálculo de probabilidades}

\ex{1} Lanzamos un dado. Calcula la probabilidad de:\begin{tasks}(3)
  \task sacar un 5
  \task sacar par
  \task sacar más de 2
  \task sacar un múltiplo de 3
  \task sacar un 7
  \task sacar menos de 7
\end{tasks}\recorda{En un experimento \textbf{regular} (resultados equiprobables): \ $P(A)=\dfrac{\text{número de casos favorables}}{\text{número de casos posibles}}$.}

\ex{1} En una bolsa hay 6 bolas rojas, 4 azules y 5 verdes. Sacamos una bola al azar. Calcula $P(\text{roja})$, $P(\text{azul})$, $P(\text{verde})$ y $P(\text{no verde})$.\par\smallskip\begin{center}\begin{tikzpicture}[x=0.42cm,y=0.42cm,baseline=(current bounding box.center)]
\draw[muted,thick,fill=soft] (-0.2,0.1) -- (-0.2,-1.70) to[out=-90,in=180] (4.00,-2.30) to[out=0,in=-90] (8.20,-1.70) -- (8.20,0.1);
\draw[muted,line width=1.6pt,line cap=round] (-0.4,0.1) -- (8.40,0.1);
\filldraw[draw=black!50,fill=red!75!black] (0.50,-0.45) circle (0.38);
\filldraw[draw=black!50,fill=red!75!black] (1.50,-0.45) circle (0.38);
\filldraw[draw=black!50,fill=red!75!black] (2.50,-0.45) circle (0.38);
\filldraw[draw=black!50,fill=red!75!black] (3.50,-0.45) circle (0.38);
\filldraw[draw=black!50,fill=red!75!black] (4.50,-0.45) circle (0.38);
\filldraw[draw=black!50,fill=red!75!black] (5.50,-0.45) circle (0.38);
\filldraw[draw=black!50,fill=blue!70!black] (6.50,-0.45) circle (0.38);
\filldraw[draw=black!50,fill=blue!70!black] (7.50,-0.45) circle (0.38);
\filldraw[draw=black!50,fill=blue!70!black] (0.50,-1.45) circle (0.38);
\filldraw[draw=black!50,fill=blue!70!black] (1.50,-1.45) circle (0.38);
\filldraw[draw=black!50,fill=teal] (2.50,-1.45) circle (0.38);
\filldraw[draw=black!50,fill=teal] (3.50,-1.45) circle (0.38);
\filldraw[draw=black!50,fill=teal] (4.50,-1.45) circle (0.38);
\filldraw[draw=black!50,fill=teal] (5.50,-1.45) circle (0.38);
\filldraw[draw=black!50,fill=teal] (6.50,-1.45) circle (0.38);
\end{tikzpicture}\end{center}

\ex{1} Lanzamos dos dados y sumamos los puntos. Completa la tabla de doble entrada y calcula la probabilidad de que la suma sea: a) 8; b) 2; c) 10 o más; d) impar.

\ex{2} Sacamos una carta de una baraja española de 40 cartas (oros, copas, espadas y bastos, con el 1 al 7, la sota, el caballo y el rey de cada palo). Calcula la probabilidad de que sea:\begin{tasks}(3)
  \task de copas
  \task un as
  \task una figura de espadas
  \task una carta que no sea figura
  \task un as o un oro
  \task el caballo de bastos
\end{tasks}

\ex{2} Escribimos cada letra de la palabra MATEMATICAS en una tarjeta y sacamos una al azar. ¿Cuál es la probabilidad de que sea una vocal? ¿Y una M? ¿Y una consonante?

\ex{2} Pablo y María juegan con un dado: Pablo gana si sale 5 o 6, y María si sale 1, 2, 3 o 4. ¿Es justo el juego? ¿Cómo lo cambiarías para que lo fuera?

\ex{2} Lanzamos dos dados. Calcula la probabilidad de que: a) saquen el mismo número; b) la diferencia entre el mayor y el menor sea 2; c) el producto sea par.

\ex{3} En una bolsa hay 5 bolas rojas y algunas azules. Si la probabilidad de sacar una bola azul es $\frac23$, ¿cuántas bolas azules hay?

\ex{3} Esta ruleta tiene tres sectores de 180°, 120° y 60°. ¿Es un experimento regular? Calcula la probabilidad de cada color.\par\smallskip\begin{center}\begin{tikzpicture}[scale=0.75,baseline=(current bounding box.center)]\fill[red!60] (0,0) -- (0:1.3) arc[start angle=0,end angle=180,radius=1.3] -- cycle;\fill[blue!55] (0,0) -- (180:1.3) arc[start angle=180,end angle=300,radius=1.3] -- cycle;\fill[teal!60] (0,0) -- (300:1.3) arc[start angle=300,end angle=360,radius=1.3] -- cycle;\draw[thick] (0,0) circle (1.3);\draw (0,0)--(0:1.3) (0,0)--(180:1.3) (0,0)--(300:1.3);\node[font=\scriptsize\bfseries] at (90:0.7) {180°};\node[font=\scriptsize\bfseries,white] at (240:0.75) {120°};\node[font=\scriptsize\bfseries] at (330:0.8) {60°};\end{tikzpicture}\end{center}


% =====================================================
\seccio{5}{Experiencias compuestas. Probabilidad}

\ex{1} Lanzamos una moneda y un dado. Haz el diagrama de árbol. ¿Cuántos resultados hay? Calcula la probabilidad de: a) cara y un 3; b) cruz y un número par; c) un número primo (con cara o cruz).\recorda{En un \textbf{diagrama de árbol}, cada camino es un resultado. Si ponemos probabilidades en las ramas, la probabilidad de un camino es el \textbf{producto} de sus ramas; si un suceso tiene varios caminos, las \textbf{sumamos}.}

\ex{1} Lanzamos dos monedas. Haz el diagrama de árbol con las probabilidades y calcula la probabilidad de sacar dos caras y la de sacar al menos una cruz.

\ex{1} La probabilidad de que llueva el sábado es 0,3 y la de que llueva el domingo, 0,6 (suponemos que son independientes). Haz el diagrama de árbol y calcula la probabilidad de que: a) llueva los dos días; b) no llueva ningún día; c) llueva al menos un día; d) llueva exactamente un día.\recorda{Si dos experiencias son \textbf{independientes}, $P(A\cap B)=P(A)\cdot P(B)$. Para «al menos uno», a menudo es más rápido el contrario: $P(\text{al menos uno})=1-P(\text{ninguno})$.}

\ex{2} Lanzamos tres monedas. Calcula la probabilidad de: a) tres caras; b) exactamente una cara; c) al menos dos caras.

\ex{2} Una jugadora de fútbol marca el 80 \% de los penaltis. Lanza dos. Haz el diagrama de árbol y calcula la probabilidad de que marque los dos, de que no marque ninguno y de que marque exactamente uno.

\ex{2} En una bolsa hay 3 bolas blancas (B) y 2 negras (N). Sacamos dos \textbf{con reemplazamiento}. Haz el diagrama de árbol y calcula la probabilidad de que: a) las dos sean blancas; b) las dos sean negras; c) salga una de cada color.\recorda{\textbf{Con reemplazamiento} devolvemos la bola a la bolsa: la segunda extracción tiene las mismas probabilidades. \textbf{Sin reemplazamiento} no la devolvemos: en la segunda extracción hay una bola menos.}

\ex{2} Lanzamos un dado tres veces. ¿Cuál es la probabilidad de no sacar ningún 6? ¿Y de sacar al menos uno?

\ex{3} Un examen tipo test tiene 3 preguntas con 4 opciones cada una (solo una correcta). Si las contestas todas al azar, ¿cuál es la probabilidad de acertarlas todas? ¿Y de no acertar ninguna? ¿Y de acertar al menos una?

\ex{3} ¿Cuántas veces tenemos que lanzar un dado, como mínimo, para que la probabilidad de sacar al menos un 6 sea mayor que $\frac12$?


% =====================================================
\seccio{6}{Experiencias dependientes. Probabilidad condicionada}

\ex{1} En una bolsa hay 5 bolas rojas y 3 verdes. Sacamos una bola, es roja y no la devolvemos a la bolsa. ¿Cuál es la probabilidad de que la segunda bola sea roja? ¿Y de que sea verde? Escríbelas con la notación de la probabilidad condicionada.\recorda{$P(B/A)$ es la probabilidad de $B$ \textbf{sabiendo que} ha pasado $A$; en el árbol, es la rama de la segunda etapa. \ $P(A\cap B)=P(A)\cdot P(B/A)$.}

\ex{1} Sacamos dos cartas de una baraja española de 40 cartas, sin reemplazamiento. a) Si la primera es un oro, ¿cuál es la probabilidad de que la segunda también lo sea? b) ¿Cuál es la probabilidad de sacar dos oros?

\ex{2} En una bolsa hay 3 bolas blancas (B) y 2 negras (N). Sacamos dos \textbf{sin reemplazamiento}. Haz el diagrama de árbol y calcula la probabilidad de que: a) las dos sean blancas; b) las dos sean negras; c) salga una de cada color. Compáralo con el ejercicio de las mismas bolas con reemplazamiento.

\ex{2} En un grupo, el 60 \% del alumnado aprueba Matemáticas. De los que aprueban Matemáticas, el 75 \% aprueba Física, y de los que no aprueban Matemáticas, solo el 25 \%. Haz el diagrama de árbol y calcula la probabilidad de que un alumno elegido al azar: a) apruebe las dos asignaturas; b) apruebe Física; c) no apruebe ninguna.

\ex{3} La urna I tiene 3 bolas rojas y 1 azul; la urna II tiene 1 bola roja y 3 azules. Lanzamos un dado: si sale 1 o 2, sacamos una bola de la urna I, y si no, de la urna II. ¿Cuál es la probabilidad de sacar una bola roja?

\ex{3} En una bolsa hay 4 bolas rojas y 2 azules. Sacamos tres, una detrás de otra, sin reemplazamiento. ¿Cuál es la probabilidad de que las tres sean rojas?

\ex{3} En una caja hay 5 bombillas y 2 son defectuosas. Probamos dos al azar, sin reemplazamiento. ¿Cuál es la probabilidad de que al menos una sea defectuosa?


% =====================================================
\seccio{7}{Tablas de contingencia}

\ex{1} En una clase de 25 alumnos hemos preguntado quién tiene mascota. Elegimos un alumno al azar. Calcula: a) $P(\text{tiene mascota})$; b) $P(\text{chico y tiene mascota})$; c) $P(\text{chica o tiene mascota})$; d) si sabemos que es una chica, la probabilidad de que tenga mascota.\par\smallskip\begin{center}{\small\renewcommand{\arraystretch}{1.3}\begin{tabular}{l|cc|c}
\rowcolor{soft} & Mascota & Sin mascota & Total\\\hline
\cellcolor{soft}\textbf{Chicas} & 8 & 4 & \textbf{12}\\
\cellcolor{soft}\textbf{Chicos} & 6 & 7 & \textbf{13}\\\hline
\cellcolor{soft}\textbf{Total} & \textbf{14} & \textbf{11} & \textbf{25}\\
\end{tabular}}\end{center}\recorda{En una \textbf{tabla de contingencia}, los casos posibles son el total (o el total de una fila, si ya sabemos en qué fila está la persona) y los favorables, las casillas que cumplen lo que piden.}

\ex{2} En un instituto hay 400 alumnos, de los cuales 180 son chicos. Hacen deporte federado 120 alumnos, 70 de los cuales son chicos. Completa la tabla y, si elegimos una persona al azar, calcula: a) $P(\text{deporte federado})$; b) $P(\text{chica y deporte federado})$; c) $P(\text{chica o deporte federado})$; d) si sabemos que es una chica, la probabilidad de que haga deporte federado.\par\smallskip\begin{center}{\small\renewcommand{\arraystretch}{1.3}\begin{tabular}{l|cc|c}
\rowcolor{soft} & Federado & No federado & Total\\\hline
\cellcolor{soft}\textbf{Chicos} & 70 & \phantom{000} & \textbf{180}\\
\cellcolor{soft}\textbf{Chicas} & \phantom{000} & \phantom{000} & \phantom{000}\\\hline
\cellcolor{soft}\textbf{Total} & \textbf{120} & \phantom{000} & \textbf{400}\\
\end{tabular}}\end{center}

\ex{2} En una empresa de 400 trabajadores, el 60 \% se ha vacunado de la gripe. De los vacunados, el 5 \% ha tenido la gripe, y de los no vacunados, el 30 \%. Completa la tabla y calcula: a) la probabilidad de que un trabajador elegido al azar haya tenido la gripe; b) $P(\text{gripe}/\text{vacunado})$ y $P(\text{gripe}/\text{no vacunado})$; c) si un trabajador ha tenido la gripe, la probabilidad de que esté vacunado.\par\smallskip\begin{center}{\small\renewcommand{\arraystretch}{1.3}\begin{tabular}{l|cc|c}
\rowcolor{soft} & Gripe & Sin gripe & Total\\\hline
\cellcolor{soft}\textbf{Vacunados} & \phantom{000} & \phantom{000} & \phantom{000}\\
\cellcolor{soft}\textbf{No vacunados} & \phantom{000} & \phantom{000} & \phantom{000}\\\hline
\cellcolor{soft}\textbf{Total} & \phantom{000} & \phantom{000} & \textbf{400}\\
\end{tabular}}\end{center}\recorda{En una tabla de contingencia, $P(A/B)=\dfrac{\text{casos de }A\text{ y }B}{\text{casos de }B}=\dfrac{P(A\cap B)}{P(B)}$: solo contamos en la fila o la columna de la condición.}

\ex{3} En un grupo de 200 personas, el 45 \% son hombres y el 30 \% lleva gafas. De las personas con gafas, 40 son mujeres. Haz la tabla de contingencia y calcula la probabilidad de que una persona elegida al azar: a) sea un hombre con gafas; b) sea una mujer sin gafas; c) sea una mujer, si sabemos que lleva gafas.

\ex{3} En una clase de 50 alumnos, 30 son de 3.\textsuperscript{o} y 20 de 4.\textsuperscript{o}. Van al instituto a pie 12 alumnos de 3.\textsuperscript{o} y 8 de 4.\textsuperscript{o}. Haz la tabla de contingencia y compara $P(\text{a pie})$ con $P(\text{a pie}/\text{3.\textsuperscript{o}})$. ¿Son independientes los sucesos «ir a pie» y «ser de 3.\textsuperscript{o}»?

\ex{3} El 70 \% de los clientes de una tienda paga con tarjeta. De los que pagan con tarjeta, el 20 \% hace una compra de más de 50 €, y de los que pagan en efectivo, el 10 \%. Haz la tabla de contingencia para 1000 clientes. Si un cliente ha hecho una compra de más de 50 €, ¿cuál es la probabilidad de que haya pagado con tarjeta?

\vfill
\begin{center}\color{muted}\small\faCheckCircle[regular]\ \ ¡Ánimo! Comprueba tus resultados con las herramientas interactivas de los apuntes de la unidad.\end{center}

\end{document}
